% ============================================================
%  style/songs.tex — Akkordzeile, Labels, Song-Umgebung, Songreihenfolge, Kategorien
%  (wird von style.tex eingebunden)
% ============================================================
% ---------- Akkordzeile & Labels ----------
\newcommand{\cl}[1]{\textbf{\color{inkchords}#1}}
\newcommand{\slbl}[1]{\vspace{0mm}\textit{\color{inkmeta}#1}}
\renewcommand{\|}{\textcolor{inkmeta}{|}}

% ---------- Automatische Nummerierung ----------
\newcounter{categorycount}
\newcounter{songcount}[categorycount]

\newcommand{\alphnum}[1]{\Alph{#1}}

% ---------- Song-Umgebung ----------
% \begin{song}{Titel}{Interpret}{Zusatzinfo (z.B. Capo)}
\newenvironment{song}[3]{%
  \stepcounter{songcount}%
  \edef\currentsongnum{\alphnum{categorycount}\thesongcount}%
  \clearpage
  \griffsongbeginn
  \thispagestyle{fancy}%
  \ifx\aktuellesongid\empty
    % Song außerhalb von \songliste: wie früher direkt ins Inhaltsverzeichnis
    \phantomsection
    \addcontentsline{toc}{section}{\makebox[\tocnumw][l]{\thesongcount}#1}%
  \else
    % Song aus \songliste: Inhaltsverzeichnis-Eintrag schreibt \songliste
    % (alphabetisch); hier nur Sprungziel, Nummer/Seite und Titel merken
    \songmarke{#1}%
  \fi
  \noindent
  \begin{minipage}[t]{20mm}
    \vspace{0pt}%
    {\songnumfont \currentsongnum}
  \end{minipage}%
  \hspace{1mm}%
  \begin{minipage}[t]{\dimexpr\linewidth-14mm\relax}
    \vspace{0pt}%
    {\songtitlefont #1\par}
    \vspace{0.8mm}
    {\songmetafont #2\ifx&#3&\else\ \textperiodcentered\ #3\fi\par}
  \end{minipage}\par
  \vspace{1.5mm}
  {\color{rulecolor}\hrule height 0.4pt}
  \vspace{1mm}%
  \setlength{\topsep}{0pt}%
  \setlength{\partopsep}{0pt}%
  \begin{alltt}\songbodysize\color{inkmain}\griffzeilenhook
}{%
  \end{alltt}
  \anmerkungausgeben
  \griffelokalzuruecksetzen
  \ifx\aktuellesongid\empty\else\songseitenmerken\fi
}

% ============================================================
%  Songreihenfolge & alphabetisches Inhaltsverzeichnis
% ============================================================
% \songliste{datei-a, datei-b, ...}   (Dateinamen aus songs/, ohne .tex)
%  - Liste ALPHABETISCH (nach Titel) führen: so erscheint sie im
%    Inhaltsverzeichnis.
%  - Im Buch werden die Songs so umsortiert, dass Songs mit gerader
%    Seitenzahl (2, 4 Seiten) immer auf einer linken (geraden) Seite
%    beginnen -> liegen als Doppelseite offen. Ein- und dreiseitige Songs
%    füllen die Lücken; sonst bleibt die Reihenfolge alphabetisch.
%    Nur wenn es gar nicht anders geht, bleibt eine Seite leer.
%  - Songnummern laufen in Buchreihenfolge.
%  - Seitenzahl je Song kommt aus dem vorigen Lauf (.aux) -> latexmk.
\makeatletter
\newcommand*{\aktuellesongid}{}
\newcommand{\songmarke}[1]{%
  % Sprungziel wie \phantomsection, aber mit festem Namen
  \edef\@currentHref{songziel.\aktuellesongid}%
  \Hy@raisedlink{\hyper@anchorstart{\@currentHref}\hyper@anchorend}%
  \protected@edef\@currentlabel{\thesongcount}%
  \label{song:\aktuellesongid}%
  \protected@write\@auxout{}{\string\songtiteldef{\aktuellesongid}{#1}}%
  \xdef\songstartseite{\the\c@page}%
}
\newcommand{\songseitenmerken}{%
  \immediate\write\@auxout{\string\songseitendef{\aktuellesongid}%
    {\the\numexpr\c@page-\songstartseite+1\relax}}%
}
\newcommand{\songtiteldef}[2]{\global\@namedef{songtitel@#1}{#2}}
\newcommand{\songseitendef}[2]{\global\@namedef{songseiten@#1}{#2}}
\newcommand{\songtitel}[1]{%
  \@ifundefined{songtitel@#1}{??}{\@nameuse{songtitel@#1}}}
\newcommand{\songseiten}[1]{% expandierbar
  \ifcsname songseiten@#1\endcsname\csname songseiten@#1\endcsname\else 1\fi}
% Inhaltsverzeichnis-Eintrag (Nummer/Seite per Label)
\newcommand{\songtoceintrag}[1]{%
  \addtocontents{toc}{\protect\contentsline{section}%
    {\protect\makebox[\protect\tocnumw][l]{\protect\ref*{song:#1}}\protect\songtitel{#1}}%
    {\protect\pageref*{song:#1}}{songziel.#1}}%
}
% leere Seite einschieben (nur wenn kein passender Song übrig ist)
\newcommand{\songleerseite}{\clearpage\thispagestyle{empty}\mbox{}}
\makeatother

\ExplSyntaxOn
\int_new:N \l__songs_i_int
\seq_new:N \l__songs_rest_seq
\prop_new:N \l__songs_file_prop
\tl_new:N \l__songs_pick_tl
\bool_new:N \l__songs_found_bool
\bool_new:N \l__songs_even_bool

\NewDocumentCommand \songliste { m }
  {
    \int_zero:N \l__songs_i_int
    \seq_clear:N \l__songs_rest_seq
    \prop_clear:N \l__songs_file_prop
    \clist_map_inline:nn {#1}
      {
        \int_incr:N \l__songs_i_int
        \tl_set:Nx \l__songs_pick_tl
          { \int_use:c { c@categorycount } - \int_use:N \l__songs_i_int }
        \seq_put_right:NV \l__songs_rest_seq \l__songs_pick_tl
        \prop_put:NVn \l__songs_file_prop \l__songs_pick_tl {##1}
        \exp_args:NV \songtoceintrag \l__songs_pick_tl
      }
    \__songs_next:
    \gdef \aktuellesongid {}
  }
\cs_new_protected:Npn \__songs_input:n #1 { \input { songs/#1 } }
\cs_new_protected:Npn \__songs_next:
  {
    \seq_if_empty:NF \l__songs_rest_seq
      {
        % beginnt der nächste Song auf einer geraden (linken) Seite?
        \int_if_even:nTF { \value{page} + 1 }
          { \bool_set_true:N \l__songs_even_bool }
          { \bool_set_false:N \l__songs_even_bool }
        \bool_set_false:N \l__songs_found_bool
        \seq_map_inline:Nn \l__songs_rest_seq
          {
            \bool_if:nT
              { \l__songs_even_bool || \int_if_odd_p:n { \songseiten {##1} } }
              {
                \tl_set:Nn \l__songs_pick_tl {##1}
                \bool_set_true:N \l__songs_found_bool
                \seq_map_break:
              }
          }
        \bool_if:NTF \l__songs_found_bool
          {
            \seq_remove_all:NV \l__songs_rest_seq \l__songs_pick_tl
            \tl_gset_eq:NN \aktuellesongid \l__songs_pick_tl
            \prop_get:NVN \l__songs_file_prop \l__songs_pick_tl \l_tmpa_tl
            \exp_args:NV \__songs_input:n \l_tmpa_tl
          }
          { \songleerseite }
        \__songs_next:
      }
  }
\ExplSyntaxOff

% 1. Zweitstimme / Hintergrundgesang
\newcommand{\voice}[1]{\textit{\color{inkmeta}#1}}

% 2. Zusatzinfos / Anweisungen
\newcommand{\info}[1]{\textit{\color{inkmeta}#1}}

% ---------- Zusammenfassung ----------
\newenvironment{summary}{%
  \vspace{-1mm}%
  \hfill%
  \begin{minipage}{0.88\linewidth}%
  \normalfont\serif\fontsize{\SzSummary}{\LdSummary}\selectfont\color{inkmain}%
  \setlength{\parindent}{0pt}%
}{%
  \end{minipage}%
  \hfill\null\par\vspace{2mm}%
}

% ---------- Kategorie-Trenner ----------
\newcommand{\category}[1]{%
  \stepcounter{categorycount}%
  \cleardoublepage
  \thispagestyle{fancy}%
  \phantomsection
  \addcontentsline{toc}{part}{\makebox[\tocnumw][l]{\alphnum{categorycount}}#1}
  \vspace*{2mm}
  {\cattitlefont #1\par}
  {\color{rulecolor}\hrule height 0.7pt}
  \vspace{4mm}
}

